\section{Introduction}\label{sec:intro}

For a primitive Dirichlet character $\chi$ of conductor $q$, let
\[
 L(s,\chi)=\sum_{n\ge1}\frac{\chi(n)}{n^s}
 \qquad(\operatorname{Re}s>1),
\]
and use its analytic continuation elsewhere. For an absolute constant
$c_1>0$, the classical region
\[
 \operatorname{Re}s\ge1-\frac{c_1}{\log(q(|\operatorname{Im}s|+2))}
\]
is zero-free except for a possible simple real zero attached to a real
character \cite[Section~14]{Davenport}. Excluding this possibility
uniformly is the Landau--Siegel zero problem. We prove its logarithmic
formulation.

\begin{theorem}\label{thm:main}
There is an absolute constant $c>0$ such that every real zero
$\beta\in(0,1)$ of every primitive nonprincipal real Dirichlet
$L$-function of conductor $q\ge3$ satisfies
\begin{equation}\label{eq:main}
 (1-\beta)\log q\ge c.
\end{equation}
\end{theorem}
